% ===========================================================================
% introduccion.tex  --  Capitulo 1
% Se incluye desde main.tex con \input{capitulos/introduccion}
% ===========================================================================


\section{Introducción}
\label{cap:introduccion}


Llevar una molécula desde su concepción hasta el mercado es un proceso lento,
caro y con una tasa de fracaso extraordinariamente alta. El recorrido completo
requiere al menos una década de trabajo y \todo{verificar y citar la cifra de
costo} DiMasi et al. (2016) reporta $\approx$2\,600 millones de USD, no 2.5
millones  \cite{mao_artificial_2026}. La ineficiencia no está distribuida de manera
uniforme: cerca del 90\,\% de las moléculas que llegan a ensayos clínicos nunca
alcanzan el registro \cite{zhang_ssgnn_2023,  czub_automlbased_2022}. Cada una de esas fallas
representa años de trabajo experimental ya consumido.
 
Parte de la dificultad es puramente combinatoria. El espacio de compuestos con
propiedades de fármaco se ha estimado en hasta $10^{60}$ estructuras, y las
bibliotecas virtuales disponibles para cribado ya superan los 36 mil millones de
compuestos \cite{hasselgren_artificial_2024}. Ninguna campaña experimental puede recorrer
ese espacio; el cribado de alto rendimiento evalúa, en el mejor de los casos,
del orden de $10^6$ moléculas. La consecuencia práctica es que la selección de
qué sintetizar y qué ensayar se toma siempre con información incompleta, y que
cualquier método capaz de ordenar candidatos antes del experimento tiene un
valor económico directo.

 
La causa de los fracasos es informativa. En la Fase II, entre el 50\,\% y el
60\,\% de los fracasos se deben a que el compuesto sencillamente no funciona,
y otro 30\,\% a señales de seguridad inaceptables; estas proporciones
prácticamente no han cambiado en dos décadas \cite{mao_artificial_2026}. Más arriba en
el embudo, aproximadamente el 40\,\% de los candidatos preclínicos se descarta
por perfiles ADMET (absorción, distribución, metabolismo, excreción y
toxicidad) insuficientes, y cerca del 30\,\% de los fármacos ya comercializados
se retira del mercado por reacciones tóxicas no previstas
\cite{zhang_computational_2025}.
 
De aquí se sigue una conclusión que organiza el resto de este trabajo: un
compuesto que se une con alta afinidad a su blanco puede ser, aun así, un
candidato inservible. La afinidad es una condición necesaria, no suficiente. Un
modelo computacional que optimice únicamente afinidad no elimina el cuello de
botella: lo desplaza a una etapa posterior, donde descartar cuesta más.
 


Este argumento no es hipotético, y el ejemplo más claro involucra precisamente
al receptor que estudia este trabajo. DSP-1181, un agonista de acción prolongada
del receptor 5-HT\textsubscript{1A} para el trastorno obsesivo-compulsivo, fue
diseñado mediante inteligencia artificial en una colaboración entre Exscientia y
Sumitomo Pharma. El diseño asistido redujo el ciclo \textit{hit-to-lead} a menos
de doce meses, un resultado notable. Sin embargo, en 2022 el compuesto se
descontinuó durante la Fase I, después de que los estudios preclínicos de
seguridad revelaran prolongación del intervalo QT en modelos animales, un factor
de riesgo conocido de arritmia cardiaca \cite{mao_artificial_2026}.
 
El caso ilustra con precisión el problema: acelerar la identificación de
\textit{hits} no garantiza la progresión clínica, y evidencia la necesidad de
integrar la predicción de seguridad y de propiedades farmacocinéticas
\emph{dentro} del ciclo de diseño, no después de él.
 

 
Los fármacos dirigidos al sistema nervioso central (SNC) enfrentan un filtro que
otros no tienen. La barrera hematoencefálica (BHE) es una barrera selectiva,
formada principalmente por el endotelio capilar cerebral y sus uniones
estrechas, que restringe el transporte paracelular de compuestos y protege al
cerebro de la exposición a sustancias potencialmente tóxicas
\cite{faramarzi_development_2022}. Un ligando con excelente afinidad por un receptor
cerebral que no cruza la BHE no tiene actividad terapéutica relevante por vía
sistémica.
 
Predecir esa permeabilidad es, en consecuencia, un obstáculo crítico en el
desarrollo de fármacos del SNC \cite{spielvogel_enhancing_2025a}. Los métodos de
referencia \textit{in vivo}, como el \textit{log\,BB} en roedor, son lentos,
costosos y requieren experimentación animal \cite{faramarzi_development_2022}. Existen
además reglas heurísticas basadas en propiedades fisicoquímicas individuales,
pero su desempeño es limitado: en una evaluación reciente sobre una base
estandarizada, un modelo de bosque aleatorio alcanzó un AUC de 0.88 para
clasificación binaria de penetración, frente a 0.53 del puntaje CNS MPO y 0.68
del BBB score \cite{spielvogel_enhancing_2025a}. La permeabilidad de la BHE es un
fenómeno multifactorial, y los modelos multivariados capturan esa estructura
mejor que cualquier descriptor aislado.
 

El receptor 5-HT\textsubscript{1A} pertenece a la familia de receptores acoplados
a proteína G y es el subtipo mejor caracterizado del sistema serotoninérgico. Se
localiza principalmente en el cerebro, en regiones mesencefálicas, límbicas y
corticales \cite{czub_automlbased_2022}. Participa en la regulación de la transmisión
neuronal, la plasticidad sináptica, la neurogénesis y la neuroprotección, y su
modulación se asocia con la reducción de ansiedad, depresión y dolor
\cite{czub_curated_2021}.
 
Su relevancia farmacológica es directa: la activación del receptor
5-HT\textsubscript{1A} es uno de los mecanismos de acción de antidepresivos,
ansiolíticos y antipsicóticos. Agonistas parciales como aripiprazol, clozapina,
buspirona, trazodona y ziprasidona, y antagonistas como risperidona, se emplean
en el tratamiento de depresión, ansiedad, esquizofrenia y trastorno bipolar
\cite{czub_automlbased_2022}. El problema es que los medicamentos disponibles presentan
numerosos efectos adversos, lo que sostiene el interés en identificar nuevas
estructuras con mejor perfil \cite{czub_automlbased_2022}. A esto se suma que la
dirección del sistema serotoninérgico es transversal: la desregulación
serotoninérgica se relaciona con migraña, dolor, ansiedad, esquizofrenia,
trastornos del movimiento y depresión, y esta última afecta a más de 300
millones de personas en el mundo \cite{czub_curated_2021}. 
 

 
La relación cuantitativa estructura--actividad (QSAR) parte de una premisa
antigua y estable: si la estructura de una molécula está codificada
adecuadamente, es posible predecir si afectará a un blanco biológico dado y con
qué intensidad \cite{czub_automlbased_2022}. La pregunta abierta no es si la estructura
contiene la información, sino \textbf{cómo representarla}.
 
El enfoque clásico fija la representación de antemano. Se calculan descriptores
fisicoquímicos y topológicos, o huellas digitales estructurales de tipo
circular, y sobre ese vector de longitud fija se entrena un modelo estadístico.
La ventaja es la transparencia y el bajo costo; la limitación es que la
representación se decide antes de ver la tarea, y por tanto no puede adaptarse a
ella.
 
Las redes neuronales de grafos (GNN) invierten ese orden. Una molécula admite una
representación natural como grafo, con los átomos como nodos y los enlaces como
aristas \cite{mao_artificial_2026}. Mediante paso de mensajes, cada nodo actualiza su
representación agregando información de sus vecinos; tras varias iteraciones,
cada átomo codifica el contexto de su vecindario químico, y una operación de
agregación global produce un vector de la molécula completa
\cite{zhang_ssgnn_2023}. La representación deja de ser un supuesto del ingeniero y
pasa a ser un parámetro aprendido junto con el resto del modelo.
 
Esa flexibilidad tiene un costo. Una GNN posee más parámetros, es más sensible a
la escasez de datos y es más difícil de auditar que un modelo de huellas
digitales con árboles de decisión. En conjuntos de tamaño moderado, la línea base
clásica frecuentemente iguala o supera a la arquitectura profunda. Por esta
razón, en este trabajo la comparación contra una línea base fuerte no es un
trámite preliminar sino un criterio de evaluación permanente
(véase el capítulo~\ref{cap:objetivos}).
 
\begin{quote}\small
\textbf{Nota sobre el lenguaje.} Todos los modelos de este trabajo son
\emph{correlacionales}. Aprenden asociaciones entre patrones estructurales y
valores medidos experimentalmente; no establecen mecanismos ni relaciones
causales. A lo largo del documento se emplea deliberadamente el verbo
\enquote{predecir} y se evitan \enquote{determinar} y \enquote{explicar}.
\end{quote}
 
\section{Planteamiento del problema}
\label{sec:intro-planteamiento}
 
Dado un conjunto de moléculas representadas por su estructura, se requiere
construir un procedimiento reproducible que estime simultáneamente tres
cantidades: (i) la afinidad de unión al receptor 5-HT\textsubscript{1A},
expresada como pKi; (ii) la probabilidad de penetración de la barrera
hematoencefálica; y (iii) el perfil de toxicidad sobre un conjunto de puntos
finales estandarizados. El objetivo del procedimiento no es sustituir la
experimentación, sino \textbf{ordenar} candidatos antes de ella.
 
El problema no es únicamente de modelado. Su dificultad principal es de calidad
de evidencia, y se concreta en cinco frentes:
 
\begin{enumerate}
  \item \textbf{Heterogeneidad y solapamiento de fuentes.} ChEMBL, BindingDB y
        PDSP contienen registros del mismo compuesto con valores distintos.
        Combinar datos de múltiples fuentes puede introducir mediciones
        contradictorias que degradan la calidad del modelo si no se resuelven
        explícitamente \cite{faramarzi_development_2022}.
  \item \textbf{Heterogeneidad de ensayo.} Ki, IC\textsubscript{50} y
        EC\textsubscript{50} no son magnitudes intercambiables; mezclarlas sin
        justificación introduce ruido sistemático en la variable objetivo.
  \item \textbf{Partición de datos.} Una división aleatoria coloca análogos
        estructurales a ambos lados de la partición y produce métricas
        optimistas que no anticipan el desempeño sobre quimiotipos nuevos.
  \item \textbf{Desbalance de clases.} Los conjuntos de toxicidad y de
        permeabilidad están fuertemente desbalanceados, situación en la que el
        ROC-AUC puede permanecer alto mientras el modelo es inútil en la clase
        minoritaria.
  \item \textbf{Dominio de aplicabilidad.} Las predicciones fuera del rango
        estructural y de respuesta del conjunto de entrenamiento son
        extrapolaciones y tienen menor probabilidad de ser confiables
        \cite{czub_automlbased_2022}.
\end{enumerate}
 
\section{Propuesta y alcance}
\label{sec:intro-propuesta}
 
Este trabajo propone un \textit{pipeline} completamente \textit{in silico} que
integra tres modelos predictivos —afinidad, permeabilidad a la BHE y
toxicidad— sobre una representación molecular común, con redes neuronales de
grafos como método principal y una comparación obligatoria contra líneas base
clásicas de huellas digitales con modelos de ensamble.
 
El alcance del presente documento corresponde al \textbf{Trabajo Terminal 1}:
adquisición y curación de los conjuntos de datos, estandarización química,
análisis exploratorio, definición del protocolo de evaluación y un modelado
inicial de referencia. La afinación de arquitecturas, la optimización de
hiperparámetros y la integración final del \textit{pipeline} corresponden al
\textbf{Trabajo Terminal 2}. Esta delimitación es deliberada: fijar el protocolo
de evaluación antes de optimizar arquitecturas evita que las decisiones de
modelado se ajusten, consciente o inconscientemente, al conjunto de prueba.
 
\section{Organización del documento}
\label{sec:intro-organizacion}
 
El capítulo~\ref{cap:objetivos} establece el objetivo general, los objetivos
específicos con sus criterios de verificación y las delimitaciones explícitas
del trabajo. El capítulo~\ref{cap:justificacion} argumenta la pertinencia del
problema, la elección del enfoque y la aportación esperada.
\todo{Completar la descripción de los capítulos restantes: marco teórico, estado
del arte, metodología, resultados preliminares y cronograma.}
 